\section{Conclusion}


Security costs performance. The costs are quantifiable:

\begin{itemize}[leftmargin=2em]
\item \textbf{Post-quantum signatures}: 51.7$\times$ larger than classical, but STARK-compressible to $O(\log N)$ bytes.
\item \textbf{FHE privacy}: $10^7$--$10^9\times$ slower than cleartext, but parallelizable and practical for targeted operations.
\item \textbf{MPC trust distribution}: 50ms--1.2s per threshold operation, acceptable for per-epoch and per-operation but not per-block.
\item \textbf{Triple-proof finality}: 248 bytes and ${\sim}2$--$5$\,ms CPU per epoch verification~\cite{lux-benchmarks}---the composition cost is small once the architecture separates time scales correctly.
\end{itemize}

The Lux architecture works because it does not ask any single primitive to do everything. BLS provides speed. ML-DSA provides post-quantum identity. Ringtail provides post-quantum privacy. STARK provides compression. MPC provides trust distribution. FHE provides data privacy. Each operates at the time scale where its cost is acceptable. The composition is the architecture.

\begin{thebibliography}{9}
\bibitem{fips204} NIST, ``FIPS 204: Module-Lattice-Based Digital Signature Standard (ML-DSA),'' August 2024.
\bibitem{cggi2020} Chillotti, Gama, Georgieva, Izabach\`{e}ne, ``TFHE: Fast Fully Homomorphic Encryption over the Torus,'' Journal of Cryptology, 2020.
\bibitem{frost2020} Komlo, Goldberg, ``FROST: Flexible Round-Optimized Schnorr Threshold Signatures,'' SAC 2020.
\bibitem{cggmp2021} Canetti, Gennaro, Goldfeder, Makriyannis, Peled, ``UC Non-Interactive, Proactive, Threshold ECDSA with Identifiable Aborts,'' CCS 2020.
\bibitem{stark} Ben-Sasson et al., ``Scalable, transparent, and post-quantum secure computational integrity,'' IACR ePrint 2018/046.
\bibitem{lux-triple} Lux Industries, Inc., ``Triple-Proof Quantum Finality: Post-Quantum Consensus for the Lux Network,'' 2026.
\bibitem{lux-tfhe} Lux Industries, Inc., ``Torus Threshold Fully Homomorphic Encryption,'' 2025.
\bibitem{lux-mpc} Lux Industries, Inc., ``M-Chain: Multi-Party Computation on the Lux Network,'' 2025.
\bibitem{lux-benchmarks} Lux Industries, ``Quasar Consensus --- Measured Benchmarks,'' 2026. Apple M1 Max, darwin/arm64, Go 1.26.1. \url{https://github.com/luxfi/consensus/blob/main/BENCHMARKS.md}.
\end{thebibliography}

